\chapter{Background}
\label{chap:background}

% 5 pages imo

In this chapter I introduce the literature background for the thesis. I begin with a comprehensive overview of existing and examined SSL training frameworks, focusing on the architectures, training, evaluation, and downstream tasks. Then, I provide a breakdown on adversarial examples, explaining the theory behind them, their universality, current attacks, and pair them with current SSL paradigm.

\section{SSL Vision Encoders}

Contrary to classification task, where the goal is to predict a correct class $c\in\{1, 2,...,K\}$ from possible $K$ classees, based on the image $x\in\mathbb{R}^{C\times H\times W}$, the goal of SSL encoders is more general. Classifiers learn based on labeled data $(x, c)$, under supervised learning paradigm,~\ie the correct prediction is well known, and the sole training objective is to teach the NN to output a specific (correct) class $c$ for each image $x$.

SSL abandons the luxury of labeled data. Intuitively, most of the data one can obtain via Internet-scale scraping is unlabeled, and the labeling requires costly human labor. Without labels, SSL posits an alternative training objective,~\ie representation learning. Representations $z\in\mathbb{R}^d$ are understood as high-dimensional ($d$ denotes number of dimensions), continuous vectors. After training, the SSL model $f_\theta: \mathbb{R}^{C\times H\times W}\rightarrow\mathbb{R}^d$ (with parameters $\theta$) should output rich and expressive representation $z$ of any image $x$, enabling further downstream use of the representations.

\subsection{Training Paradigms}

Notably, what exactly is a "rich and expressive representation $z$ of any image $x$, enabling further downstream use of the representations" is not well-grounded in theory. Moreover, one cannot compute gradient of such an objective. 

More precise view of the above goal is: representations $z_1$, $z_2$ of two \textit{similar} images $x_1$, $x_2$ should be \textit{closer} than representations $z_1$, $z_3$ of two \textit{dissimilar} images $x_1$, $x_3$. Intuitively, an image of a dog should have a representation that lies close to other images of dogs in the embedding space, and should have a representation that lies far from images depicting circuit boards. 

\subsection{Similar Images and SSL}
\label{sec:similar_imgs_ssl}

Defining "similarity" of images is tricky, especially since images can be similar in many ways. The similarity may be structural, compositional, semantic, stylistic, etc. To avoid ambiguities and human-based biases, the default approach is to augment \textit{the same} image multiple times, producing multiple "views". These views are (likely) similar to each other, and the similarity is as broad as imaginable. 

Representations' "closeness" is easier to precisely formulate. Since embeddings are continuous vectors, the natural approach is to use an abstract \textit{distance} function $\text{dist}: \mathbb{R}^d\times\mathbb{R}^d\rightarrow\mathbb{R}$, for example L2 distance or cosine similarity. 

Effectively, given a dataset of $N$ unlabeled images $D=\{x_i\},i=1,...,N$ the SSL training objective utilizing multiple views of the same image is as follows:

\begin{equation}
    \label{eq:loss_ssl_sim}
    \theta^*=\underset{\theta}{\argmin}
        \mathbb{E}_{x\sim D,\mathcal{A}_i\sim{\mathbb{A}},\mathcal{A}_j\sim\mathbb{A}}
        \left[
        \text{dist}(\mathcal{A}_i(x),\mathcal{A}_j(x))
        \right]\text{,}
\end{equation}
where $\mathbb{A}=\{\mathcal{A}_1,\mathcal{A}_2,...\}$ is a set of augmentations $\mathcal{A}_i: \mathbb{R}^{C\times H\times W} \rightarrow \mathbb{R}^{C\times H\times W}$.

\subsection{The Great Collapse}

Unfortunately,~\cref{eq:loss_ssl_sim} has a fatal flaw: its minimizer is a network $f_\theta^*: \mathbb{R}^{C\times H\times W}->[0]^d$. To address that HERE SIMSIAM AND DINO ONCE I HAVE INTERNET TO WRITE IT AS GOOD AS IT GETS.

\subsection{Dissimilar Images and Contrastive Learning}

While \textit{similar} images' embeddings should stay close, it is also important that \textit{dissimilar} images' embeddings stay far from each other. Analogously to~\cref{sec:similar_imgs_ssl}, it is not clear what does exactly mean that two images are dissimilar. 

The trivial solution is to assume that two different images are dissimilar from the definition. Contrastive Learning assumes exactly that, with the InfoNCE training loss defined as:

\begin{equation}
    \label{eq:infonce}
    \text{InfoNCE}=
\end{equation}

Intuitively, this training objective forces the NN to "push away" the representations of all alterantive images from the original positive image. Such a paradigm has been widely adopted HERE MORE ABOUT SIMCLR, AND MAYBE ABOUT CLIP.

ALSO ABOUT THE TRAINING DATA, IMAGENET, HOW MANY EPOCHS, THAT A LOT OF STUFF IS STANDARDIZED ETC.

\subsection{Evaluating SSL Image Encoders}

The generality of the SSL training task (\cref{eq:loss_ssl_sim,eq:infonce}) comes at a cost. One can not simply look at the produced embeddings $z$ and assess their quality, since their dimension $d$ is often too high (above 100, sometimes even 1000) to allow for meaningful analysis, yet alone a universal quality metric.

Instead, researchers focus on measuring how useful the embeddings are to help solve well-defined tasks, like classification or semantic segmentation. For classification, the dominant approach is to select a downstream dataset $D_T=\{(x_i,c_i)\}, i=1,...,N_T$ of $N_T$ labeled images, and based on the representations from trained $f_\theta$, fit a simple linear model $l_\mathbf{w}: \mathbb{R}^d \rightarrow \{1,...,K\}$ minimizing cross-entropy loss $\text{CELoss}=-\log(1-p_i)$, where $p_i$ is the probability of i-th class returned by $l_\mathbf{w}$. The final encoder+linear model classification stack is defined as

\begin{equation}
    (l_\mathbf{w}\circ f_\theta): \mathbb{R}^{C\times H\times W} \rightarrow \{1,...,K\}\text{,}
\end{equation}
and its downstream task performance is evaluated as accuracy on the validation set of the training data $D_T$. The established approach is to perform a single training epoch on the downstream data.

TODO HERE ABOUT segmentation


\section{Adversarial Examples}

While most of the field focused on making the models better and more general, TODO CITET indentified a grave vulnerability of almost every neural network that exists. With some assumptions (introduced later) and little effort, any NN can be forced to output a wrong prediction with high consistency. What is even worse, such an adversarial input to the NN can be made indistinguishable from a real, benign input, and can evade detection. Changing a few pixel values by extremely small values (for example by at most 8/255 of the full pixel range) is enough to reliably fool a vast majority of image-based NNs, regardless of their training data, training objective, architecture, hyperparameters, preprocessing, and so on.

TODO CITIT AGAIN dubbed their discovery as "adversarial examples". More precisely, an adversarial example $x_{adv}=x+\delta$ is a benign sample $x$ (here: image) and its perturbation $\delta$, crafted explicitly to cause misprediction of a victim model $f_V$. To limit the difference between clean and adversarial image, $\delta$ is oftentimes constrained to be $\left\|\delta\right\|_p<\epsilon$, where $p\in\{0,1,2,\infty\}$. To ensure imperceptibility, $\epsilon$ is usually very small,~\eg $8/255$ for $L_\infty$ norm, meaning each pixel of the image is changed by at most $8/255$ (assuming pixel values are in $[0,1]$).

Another important element of adversarial examples is the optimization procedure to obtain $\delta$. Given an adversarial objective, 

Here more about the theory that the errors will accumulate etc

PGD here

SegPGD here































% This chapter summarizes the technical background needed for the robustness evaluation: reusable vision encoders, SSL objectives, downstream tasks, adversarial attacks, and robust SSL defenses.

% \section{Foundation Vision Encoders}

% Foundation models are trained at scale and then adapted to tasks that are not fixed during pre-training \parencite{bommasani2021opportunities}. In vision, the foundation component is often an encoder $f_\theta$ that maps an image $x$ to a representation $z=f_\theta(x)$. A downstream adaptor $h_\phi$ maps $z$ to the task output. For classification the output is a class label; for segmentation it is a pixel-label map; for depth estimation it is a dense continuous depth map.

% Vision Transformers (ViTs) are a natural backbone for this setting \parencite{dosovitskiy2020image}. An image is divided into patches, each patch is embedded as a token, and transformer blocks process the token sequence. Classification can use a class token or pooled token, while dense tasks use patch tokens and upsampling heads. The architecture therefore exposes both global and spatial representations.

% \section{Self-Supervised Learning}

% SSL learns representations without manual labels. A common formulation samples two augmentations $x_1$ and $x_2$ of the same image and trains the encoder so that their representations agree \parencite{chen2020simple,balestriero2023cookbook}. SimCLR uses a contrastive objective that attracts positive pairs and repels negatives \parencite{chen2020simple}. SimSiam removes explicit negatives and uses a stop-gradient Siamese structure \parencite{simsiam}.

% \DINO uses self-distillation \parencite{caron2021dino}. A student network predicts the output distribution of a teacher network on a different augmented view; the teacher is updated as an exponential moving average of the student. Centering and sharpening avoid collapse. \DINO features are especially relevant because they transfer well to ViT downstream tasks.

% \DINOv extends this line at larger scale and with stronger training recipes \parencite{oquab2024dinov2}. It is important for this thesis because its representations are explicitly used for diverse tasks, including dense prediction. A robustness evaluation that only tests linear classification is therefore narrower than the claimed deployment value.

% \section{Downstream Tasks}

% \textbf{Linear classification} is the standard SSL evaluation protocol. The encoder is frozen and a linear classifier is trained on top of its representation \parencite{chen2020simple,simsiam}. Accuracy measures how linearly separable the learned representation is.

% \textbf{Semantic segmentation} assigns a class to every pixel. Fully convolutional networks established the dense prediction setting \parencite{long2017segmentation}, and UPerNet introduced unified multi-scale parsing for scene understanding \parencite{xiao2018unified}. For frozen ViTs, segmentation heads consume patch tokens, upsample logits, and optimize pixel-wise cross-entropy. Performance is reported as mean intersection-over-union (mIoU).

% \textbf{Monocular depth estimation} predicts a depth value from a single RGB image. It is dense like segmentation, but the output is continuous and geometric. Depth training often combines pixel-wise scale-invariant errors with gradient or smoothness terms \parencite{godard2019depthestimation,MegaDepthLi18,AdaBins}. This thesis reports root mean squared error (RMSE), where lower is better.

% \section{Adversarial Robustness}

% Adversarial examples are small perturbations that cause incorrect or degraded predictions \parencite{biggio2013evasion,szegedy2014intriguing}. For a task loss $L$, a white-box $\ell_\infty$ attack solves
% \begin{equation}
%     \delta^\star =
%     \argmax_{\lVert\delta\rVert_\infty \leq \varepsilon}
%     L(h_\phi(f_\theta(x+\delta)), y).
% \end{equation}
% Projected gradient descent (PGD) approximately solves this objective through iterative gradient ascent and projection back to the allowed perturbation ball \parencite{madry2018towards}.

% For SSL encoders, the attacked object is ambiguous. Without a downstream head, an attacker can maximize representation distance. With a deployed head, the attacker can maximize classification cross-entropy, segmentation loss, or depth loss. The distinction is central: a defense optimized for one objective may not protect the others.

% \section{Robust Self-Supervised Learning}

% Robust SSL methods adapt adversarial training to unlabeled representation learning. RoCL aligns clean and adversarial contrastive views \parencite{kim2020adversarial}. ACL enforces adversarial consistency during contrastive learning \parencite{jiang2020robust}. AdvCL incorporates pseudo-label information \parencite{fan2021does}. CLAE constructs adversarial contrastive examples \parencite{ho2020CLAE}. DYNACL studies the tension between strong augmentations and robustness, using dynamic augmentation schedules \parencite{luo2023rethinking}.

% \DeACL is the defense evaluated here \parencite{zhang2022decoupled}. Instead of robustly pre-training from scratch, it starts from a pre-trained encoder $f$ and fine-tunes a robust encoder $f_R$. The objective preserves clean representations while aligning clean and adversarial representations:
% \begin{equation}
%     L(f_R,f) = d(f_R(x), f(x)) + \gamma d(f_R(x_{\mathrm{adv}}), f_R(x)).
%     \label{eq:background-deacl}
% \end{equation}
% The adversarial example $x_{\mathrm{adv}}$ is produced by maximizing a representation-distance objective. This design makes \DeACL computationally attractive, but it also motivates the thesis question: a defense trained on representation distance may not transfer to downstream losses.
